\documentclass[a4paper,fontset=windows]{ctexart}

\usepackage{anyfontsize}
\usepackage{amsmath,amssymb,mathrsfs,bm,wasysym,esint}
\allowdisplaybreaks
\usepackage{indentfirst}
\setlength{\parindent}{0em}
\usepackage{geometry}
\geometry{top=3cm,bottom=3cm,left=2cm,right=2cm}

\begin{document}

\title{\textbf{数学物理方法作业}}
\author{1900011604\ \ 张植竣\ \ 北京大学}
\date{2020年12月5日}
\maketitle

\section*{第十八章}
\bigskip

\begin{itemize}

\item[1.(1)]
将$x\dfrac{\mathrm{d}^2 y}{\mathrm{d}x^2} + 2\dfrac{\mathrm{d}y}{\mathrm{d}x} + (\lambda+x)y = 0$与标准形式：
\[
\frac{\mathrm{d}}{\mathrm{d}x}\left[p(x)\frac{\mathrm{d}y}{\mathrm{d}x}\right] + \left[\lambda\rho(x) - q(x)\right]y = p(x)\frac{\mathrm{d}^2 y}{\mathrm{d}x^2} + p'(x)\frac{\mathrm{d}y}{\mathrm{d}x} + \left[\lambda\rho(x) - q(x)\right]y = 0
\]
对比得：
\[
\frac{p(x)}{p'(x)} = \frac{x}{2} \quad\Rightarrow\quad \frac{\mathrm{d}p}{p} = \frac{2\mathrm{d}x}{x} \quad\Rightarrow\quad p(x) = Cx^2
\]
取$p(x) = x^2$得：
\[
x^2\frac{\mathrm{d}^2 y}{\mathrm{d}x^2} + 2x\frac{\mathrm{d}y}{\mathrm{d}x} + \left[\lambda\rho(x) - q(x)\right]y
\]
原式等号两边同乘$x$得：
\[
x^2\dfrac{\mathrm{d}^2 y}{\mathrm{d}x^2} + 2x\dfrac{\mathrm{d}y}{\mathrm{d}x} + (\lambda x+x^2)y = 0
\]
对比得：
\[
\rho(x) = x,\qquad q(x) = -x^2
\]
化为标准形式：
\[
\frac{\mathrm{d}}{\mathrm{d}x}\left(x^2\frac{\mathrm{d}y}{\mathrm{d}x}\right) + \left(\lambda x+x^2\right)y = 0
\]
\medskip

\item[1.(3)]
将$x(1-x)\dfrac{\mathrm{d}^2 y}{\mathrm{d}x^2} + (1-x)\dfrac{\mathrm{d}y}{\mathrm{d}x} + \lambda y = 0$与标准形式：
\[
\frac{\mathrm{d}}{\mathrm{d}x}\left[p(x)\frac{\mathrm{d}y}{\mathrm{d}x}\right] + \left[\lambda\rho(x) - q(x)\right]y = p(x)\frac{\mathrm{d}^2 y}{\mathrm{d}x^2} + p'(x)\frac{\mathrm{d}y}{\mathrm{d}x} + \left[\lambda\rho(x) - q(x)\right]y = 0
\]
对比得：
\[
\frac{p(x)}{p'(x)} = \frac{x}{1-x} \quad\Rightarrow\quad \frac{\mathrm{d}p}{p} = \frac{(1-x)\mathrm{d}x}{x} \quad\Rightarrow\quad p(x) = Cx\mathrm{e}^{-x}
\]
取$p(x) = x\mathrm{e}^{-x}$得：
\[
x\mathrm{e}^{-x}\frac{\mathrm{d}^2 y}{\mathrm{d}x^2} + (1-x)\mathrm{e}^{-x}\frac{\mathrm{d}y}{\mathrm{d}x} + \left[\lambda\rho(x) - q(x)\right]y
\]
原式等号两边同乘$\mathrm{e}^{-x}$得：
\[
x\mathrm{e}^{-x}\dfrac{\mathrm{d}^2 y}{\mathrm{d}x^2} + (1-x)\mathrm{e}^{-x}\dfrac{\mathrm{d}y}{\mathrm{d}x} + \lambda y = 0
\]
对比得：
\[
\rho(x) = \mathrm{e}^{-x},\qquad q(x) = 0
\]
化为标准形式：
\[
\frac{\mathrm{d}}{\mathrm{d}x}\left(x\mathrm{e}^{-x}\frac{\mathrm{d}y}{\mathrm{d}x}\right) + \lambda\mathrm{e}^{-x}y = 0
\]
\bigskip

\item[2.]
原式等号两边同乘$r^2$得：
\[
r\frac{\mathrm{d}}{\mathrm{d}r}\left(r\frac{\mathrm{d}R}{\mathrm{d}r}\right) + \lambda R = r^2\frac{\mathrm{d}^2 R}{\mathrm{d}r^2} + r\frac{\mathrm{d}R}{\mathrm{d}r} + \lambda R = 0
\]
设$R = r^\rho$得：
\[
\rho(\rho-1)r^{\rho-2} \cdot r^2 + \rho r^{\rho-1} \cdot r + \lambda r^\rho = 0 \quad\Rightarrow\quad (\rho^2+\lambda)r^\rho = 0
\]
即$\rho^2+\lambda = 0$，分以下两种情况：

\begin{itemize}

\item[$\bullet$]
$\lambda = -m^2\ (m\in\mathbb{R}^+)$，解得$\rho = \pm m$。

则$R = Ar^\rho + Br^{-\rho} = Ar^m + Br^{-m}$。

代入边界条件得：
\[
\left\{
\begin{aligned}
&R(a) = Aa^m + Ba^{-m} = 0 \\
&R(b) = Ab^m + Bb^{-m} = 0
\end{aligned}
\right.
\]
这是一个关于$A$和$B$的二元线性方程组，其解为：
\[
A = \frac{\begin{vmatrix}0 & a^{-m} \\ 0 & b^{-m}\end{vmatrix}}{\begin{vmatrix}a^m & a^{-m} \\ b^m & b^{-m}\end{vmatrix}} = 0,\qquad B = \frac{\begin{vmatrix}a^m & 0 \\ b^m & 0\end{vmatrix}}{\begin{vmatrix}a^m & a^{-m} \\ b^m & b^{-m}\end{vmatrix}} = 0
\]
当且仅当$\begin{vmatrix}a^m & a^{-m} \\ b^m & b^{-m}\end{vmatrix} = 0$，即$a = \pm b$时，方程有非零解，但此时方程解不唯一，且与已知条件$b>a>0$不符。

故$\lambda = -m^2\ (m\in\mathbb{R}^+)$不合题意，舍去。

\item[$\bullet$]
$\lambda = m^2\ (m\in\mathbb{R}^+)$，解得$\rho = \pm im$。

则$R = Ar^{im} + Br^{-im} = A\mathrm{e}^{im\ln{r}} + B\mathrm{e}^{-im\ln{r}} = C\cos{m\ln r} + D\sin{m\ln r}$。

考虑边界条件$R(a) = R(b) = 0$，不便于确定本征值$\lambda$，应对$r$换元。

设$z = \dfrac{r}{a}$，则$r = az$，$\dfrac{\mathrm{d}z}{\mathrm{d}r} = \dfrac{1}{a}$，得：
\[
\frac{1}{r}\frac{\mathrm{d}}{\mathrm{d}r}\left(r\frac{\mathrm{d}R}{\mathrm{d}r}\right) + \frac{\lambda}{r^2}R = \frac{1}{az}\cdot\frac{\mathrm{d}z}{\mathrm{d}r}\cdot\frac{\mathrm{d}}{\mathrm{d}z}\left(az\cdot\frac{\mathrm{d}z}{\mathrm{d}r}\cdot\frac{\mathrm{d}R}{\mathrm{d}z}\right) + \frac{\lambda}{a^2 z^2}R = 0
\]
\[
\frac{1}{a^2}\left[\frac{1}{z}\frac{\mathrm{d}}{\mathrm{d}z}\left(z\frac{\mathrm{d}R}{\mathrm{d}z}\right) + \frac{\lambda}{r^2}R\right] = 0
\]
换元后方程形式与原方程相同，其解为：$R = M\cos{m\ln z} + N\sin{m\ln z}$。

边界条件变为：
\[
R(z = 1) = R(z = \dfrac{b}{a}) = 0
\]
此时$\ln1 = 0$，仅有一个关于$m$的方程，便于确定本征值。

代入边界条件得：
\[
\left\{
\begin{aligned}
&M = 0 \\
&M\cos{m\ln\frac{b}{a}} + N\sin{m\ln\frac{b}{a}} = 0
\end{aligned}
\right.
\]
得：
\[
m\ln\frac{b}{a} = n\pi \quad\Rightarrow\quad m = \frac{n\pi}{\ln b - \ln a}\qquad (n\in\mathbb{N}^+)
\]
\end{itemize}

由$\lambda = m^2$得本征值：
\[
\lambda = \left(\frac{n\pi}{\ln b - \ln a}\right)^2
\]
取$N = 1$得本征函数：
\[
R = \sin{m\ln\frac{r}{a}} = \sin{\frac{\ln r - \ln a}{\ln b - \ln a}n\pi}
\]
综上，本征值问题的解为：
\[
\lambda = \left(\frac{n\pi}{\ln b - \ln a}\right)^2
\]
\[
R = \sin{m\ln\frac{r}{a}} = \sin{\frac{\ln r - \ln a}{\ln b - \ln a}n\pi}
\]
\bigskip

\item[3.]
设对应于$\lambda_1$、$\lambda_2$的本征函数分别为$y_1(x)$、$y_2(x)$，则对应于$\lambda_2^*$的本征函数为$y_2^*(x)$。

则有方程组：
\begin{equation}
\frac{\mathrm{d}}{\mathrm{d}x}\left[p\frac{\mathrm{d}y_1}{\mathrm{d}x}\right] + \left[\lambda_1\rho - q\right]y_1 = 0
\end{equation}
\begin{equation}
\frac{\mathrm{d}}{\mathrm{d}x}\left[p\frac{\mathrm{d}y_2^*}{\mathrm{d}x}\right] + \left[\lambda_2^*\rho - q\right]y_2^* = 0
\end{equation}
$(1)\times y_2^* - (2)\times y_1$得：
\[
y_2^*\frac{\mathrm{d}}{\mathrm{d}x}\left[p\frac{\mathrm{d}y_1}{\mathrm{d}x}\right] + \left[\lambda_1\rho - q\right]y_1y_2^* - y_1\frac{\mathrm{d}}{\mathrm{d}x}\left[p\frac{\mathrm{d}y_2^*}{\mathrm{d}x}\right] + \left[\lambda_2^*\rho - q\right]y_2^*y_1 = 0
\]
积分得：
\[
\int_{a}^{b}(\lambda_2^* - \lambda_1)\rho(x)y_1(x)y_2^*(x)\,\mathrm{d}x = \left.p(x)\left[y_2^*y'_1 - {y_2^*}'y_1\right]\right|_a^b
\]
又因为$p(a) = p(b)$，则有：
\[
\int_{a}^{b}(\lambda_2^* - \lambda_1)\rho(x)y_1(x)y_2^*(x)\,\mathrm{d}x = 0 \quad\Leftrightarrow\quad \left.\left[y_2^*y'_1 - {y_2^*}'y_1\right]\right|_a^b = 0
\]
即：
\[
\int_{a}^{b}(\lambda_2^* - \lambda_1)\rho(x)y_1(x)y_2^*(x)\,\mathrm{d}x = 0 \quad\Leftrightarrow\quad y_2^*(b)y'_1(b) - {y_2^*}'(b)y_1(b) = y_2^*(a)y'_1(a) - {y_2^*}'(a)y_1(a)
\]
由题意得：（注意${y_2^*}' = {y'_2}^*$）
\begin{align*}
&\quad\ \, y_2^*(b)y'_1(b) - {y_2^*}'(b)y_1(b) \\
&= \left[\alpha_{11}y_2^*(a) + \alpha_{12}{y'_2}^*(a)\right]\left[\alpha_{21}y_1(a) + \alpha_{22}y'_1(a)\right] - \left[\alpha_{21}y_2^*(a) + \alpha_{22}{y'_2}^*(a)\right]\left[\alpha_{11}y_1(a) + \alpha_{12}y'_1(a)\right] \\
&= \begin{vmatrix}\alpha_{11} & \alpha_{11} \\ \alpha_{21} & \alpha_{21}\end{vmatrix}y_2^*(a)y_1(a) + \begin{vmatrix}\alpha_{11} & \alpha_{11} \\ \alpha_{22} & \alpha_{21}\end{vmatrix}{y_2^*}'(a)y_1(a) + \begin{vmatrix}\alpha_{11} & \alpha_{12} \\ \alpha_{21} & \alpha_{22}\end{vmatrix}y_2^*(a)y'_1(a) + \begin{vmatrix}\alpha_{12} & \alpha_{12} \\ \alpha_{22} & \alpha_{22}\end{vmatrix}{y_2^*}'(a)y'_1(a)\\
&= \begin{vmatrix}\alpha_{11} & \alpha_{12} \\ \alpha_{21} & \alpha_{22}\end{vmatrix}\left[y_2^*(a)y'_1(a) - {y_2^*}'(a)y_1(a)\right]
\end{align*}
故$\begin{vmatrix}\alpha_{11} & \alpha_{12} \\ \alpha_{21} & \alpha_{22}\end{vmatrix} = 1$时：
\[
y_2^*(b)y'_1(b) -{y_2^*}'(b)y_1(b) = y_2^*(a)y'_1(a) - {y_2^*}'(a)y_1(a)
\]
即：
\[
\int\limits_{a}^{b}(\lambda_2^* - \lambda_1)\rho(x)y_1(x)y_2^*(x)\,\mathrm{d}x = 0
\]
本征函数正交性得证。

\end{itemize}

\end{document}
